\section{Definitions and certificate principles}
\label{sec:preliminaries}

All graphs in this manuscript are finite and undirected.  A Hamiltonian cycle
is regarded as an undirected edge set; reversal and cyclic rotation do not
produce new cycles.

\begin{definition}[Barnette graph]
A \emph{Barnette graph} is a simple, planar, bipartite, cubic, and
3-vertex-connected graph.
\end{definition}

Let \(\Hcal(G)\) denote the set of undirected Hamiltonian cycles of a
Hamiltonian graph \(G\).

\begin{definition}[Hamiltonian edge separation]
For distinct edges \(e,f\in E(G)\), a Hamiltonian cycle \(C\) \emph{covers} the
ordered requirement \((e,f)\) if \(e\in C\) and \(f\notin C\).  A family
\(\Ccal\subseteq\Hcal(G)\) is \emph{Hamiltonian edge-separating} if it covers
every ordered pair of distinct edges.  Define
\begin{equation}
\label{eq:hsep-definition}
 \hsep(G)=
 \min\left\{ |\Ccal|:
 \begin{array}{l}
 \Ccal\subseteq\Hcal(G),\\[-2pt]
 \text{for all distinct }e,f\in E(G),
 \text{ some }C\in\Ccal\text{ has }e\in C, f\notin C
 \end{array}
 \right\}.
\end{equation}
\end{definition}

\begin{definition}[Barnette extremal sequence]
For an order at which the relevant values are defined, put
\begin{equation}
\label{eq:MB-definition}
 \MB(n)=\max\{\hsep(G):G\text{ is a Barnette graph and }|V(G)|=n\}.
\end{equation}
If no Barnette graph of order \(n\) exists, \(\MB(n)\) is undefined.  In the
certified range every graph in every nonempty census is Hamiltonian, so
\cref{eq:MB-definition} is unambiguous there.
\end{definition}

Given an ordered family \(C_1,\ldots,C_k\), associate with an edge \(e\) its
incidence signature
\[
 \sigma(e)=\{i:e\in C_i\}\subseteq\{1,\ldots,k\}.
\]

\begin{proposition}[Signature antichain]
\label{prop:signature-antichain}
A family \(C_1,\ldots,C_k\) is Hamiltonian edge-separating if and only if the
edge signatures form an antichain: for all distinct \(e,f\), neither
\(\sigma(e)\subseteq\sigma(f)\) nor \(\sigma(f)\subseteq\sigma(e)\).
\end{proposition}

\begin{proof}
The requirement \((e,f)\) is covered exactly when some index belongs to
\(\sigma(e)\setminus\sigma(f)\).  Covering both ordered requirements
\((e,f)\) and \((f,e)\) is therefore equivalent to incomparability of the two
signatures.
\end{proof}

In particular, if \(G\) has \(m\) edges and \(k\) selected cycles, Sperner's
theorem gives the necessary inequality
\[
 m\le \binom{k}{\lfloor k/2\rfloor}.
\]
This general counting obstruction is kept separate from graph-specific
optimality certificates.

\begin{lemma}[Primal and packing certificates]
\label{lem:certificate-principle}
If \(k\) explicit Hamiltonian cycles cover every ordered edge-pair requirement,
then \(\hsep(G)\le k\).  If there are \(k\) requirements such that every
Hamiltonian cycle covers at most one of them, then \(\hsep(G)\ge k\).
Consequently, matching certificates prove \(\hsep(G)=k\).
\end{lemma}

\begin{proof}
The first statement is the definition.  For the second, each of the \(k\)
requirements must be assigned to a selected cycle that covers it, while no
selected cycle can receive two requirements.  Hence at least \(k\) cycles are
necessary.
\end{proof}

When no matching packing is available, a lower bound in this project is called
independently certified only if a standalone verifier checks a complete
Hamiltonian-cycle universe and exhaustively excludes every cover below the
claimed value.
